% !TEX program = pdflatex
\documentclass[aspectratio=169,11pt]{beamer}

\usetheme{Madrid}
\usecolortheme{seahorse}
\setbeamertemplate{navigation symbols}{}
\setbeamertemplate{footline}[frame number]

\usepackage[T1]{fontenc}
\usepackage{lmodern}
\usepackage{listings}
\usepackage{booktabs}
\usepackage{tikz}
\usetikzlibrary{arrows.meta,positioning}

\lstset{
  basicstyle=\ttfamily\scriptsize,
  keywordstyle=\color{blue!70!black}\bfseries,
  commentstyle=\color{green!50!black},
  stringstyle=\color{red!60!black},
  breaklines=true,
  showstringspaces=false,
  columns=fullflexible,
  frame=single,
  backgroundcolor=\color{gray!8},
  rulecolor=\color{gray!40},
  xleftmargin=2pt,
  xrightmargin=2pt,
  language=Java,
}

\definecolor{beastblue}{RGB}{40,80,140}
\definecolor{beastorange}{RGB}{200,100,20}
\setbeamercolor{frametitle}{fg=beastblue}
\setbeamercolor{title}{fg=beastblue}
\setbeamercolor{structure}{fg=beastblue}

\newcommand{\pkg}[1]{\texttt{\color{beastorange}#1}}
\newcommand{\code}[1]{\texttt{#1}}

\title{The Scalable Contract}
\subtitle{Resolving beast3 issue \#20: a binding three-method contract on \texttt{Scalable}}
\author{Alexei Drummond}
\date{5 May 2026}

\begin{document}

% ============================================================
\begin{frame}
\titlepage
\end{frame}

% ============================================================
\begin{frame}{The blocker}
Joelle Barido-Sottani has been blocked porting custom-tree packages to beast3 since 26 March.

\vspace{1em}

Issue \href{https://github.com/CompEvol/beast3/issues/20}{CompEvol/beast3\#20}, opened by Remco in October 2025.

\vspace{1em}

Two distinct problems she ran into:

\begin{enumerate}
\item The old \code{Scalable} interface returns degrees of freedom (\code{int}), which encodes a specific Hastings-ratio shape (\code{dof $\times$ log(s)}).
For her custom trees the HR isn't of that shape.

\item The new spec \code{UpDownOperator} bypasses \code{Tree.scale()} entirely and runs its own recursion.
Even if she fixed her tree's \code{scale()}, UpDown wouldn't call it.
\end{enumerate}
\end{frame}

% ============================================================
\begin{frame}[fragile]{The old contract}
\begin{lstlisting}
public interface Scalable {
    int  scale(double s);            // returns dof
    void scaleOne(int i, double s);
    default double scaleAll(double s) {
        return scale(s) * Math.log(s);
    }
}
\end{lstlisting}

\vspace{0.5em}

\textbf{What's wrong:}

\begin{itemize}
\item \code{int dof} is a stand-in for one specific HR formula.
Operators reconstruct \code{dof $\times$ log(s)} as the move's Jacobian contribution.
\item Anything that scales differently (interval scaling, custom trees, IntervalScaleOperator) has to bypass the interface and compute HR inline.
\item The interface gives no way to ask ``what does this Scalable use as its dilation axis?''
\end{itemize}
\end{frame}

% ============================================================
\begin{frame}{What does \texttt{scale} actually return?}
The Metropolis-Hastings acceptance ratio has three parts:

\begin{enumerate}
\item \textbf{Target ratio}: $\pi(x') / \pi(x)$.
\item \textbf{Proposal ratio} (the Hastings ratio proper): $q(x \mid x') / q(x' \mid x)$.
\item \textbf{Jacobian determinant}: $|\det(\partial x' / \partial x)|$ from the change of variables.
\end{enumerate}

\vspace{1em}

\code{Scalable.scale(s)} can only sensibly return the \textbf{Jacobian} of the deterministic part of the move.

\vspace{0.5em}

The proposal ratio depends on how the operator drew \code{s} (uniform-on-log, Bactrian, Gaussian, ...) and belongs to the operator's kernel.

\vspace{0.5em}

\textbf{Joelle's option (1)} of returning ``the HR'' is right in spirit; the precise refinement is to return the log Jacobian determinant.
\end{frame}

% ============================================================
\begin{frame}[fragile]{The new contract}
\begin{lstlisting}
public interface Scalable {
    double scale(double s);            // returns log Jacobian
    double getScalableValue();         // current pos. on dilation axis
    default double setScalableValue(double V) {
        return scale(V / getScalableValue());
    }
    void scaleOne(int i, double s);
}
\end{lstlisting}

\vspace{0.5em}

Three methods, mutually consistent on a single dilation axis.
\end{frame}

% ============================================================
\begin{frame}[fragile]{The three invariants}
For any valid Scalable \code{x} and positive \code{s}:

\begin{enumerate}
\item \textbf{Scale-equivariance.} After \code{scale(s)}, \code{getScalableValue()} returns $s \times$ its previous value.
\begin{lstlisting}
double v0 = x.getScalableValue();
x.scale(s);
assert x.getScalableValue() == s * v0;
\end{lstlisting}

\item \textbf{Set is a fixed point of get.}
\begin{lstlisting}
x.setScalableValue(V);
assert x.getScalableValue() == V;
\end{lstlisting}

\item \textbf{Set composes with scale.}
\begin{lstlisting}
x.setScalableValue(x.getScalableValue() * s) == x.scale(s);
\end{lstlisting}
\end{enumerate}

\vspace{0.5em}

The default \code{setScalableValue} expresses (3) directly.
\end{frame}

% ============================================================
\begin{frame}{What is the dilation axis?}
\code{getScalableValue} must return whatever quantity \emph{actually} scales by exactly \code{s} under the implementer's \code{scale}.

\vspace{1em}

\begin{tabular}{ll}
\toprule
\textbf{Implementer} & \textbf{Dilation summary} \\
\midrule
\code{RealScalarParam} & the value itself \\
\code{RealVectorParam} & sum of values \\
\code{Tree} (interval scaling) & sum of margins \\
\code{Tree} (affine height scaling) & root height \\
Joelle's custom tree & whatever her scale preserves \\
\bottomrule
\end{tabular}

\vspace{1em}

\textbf{The contract is enforced per-implementation.}
A type that violates the invariants is not a valid Scalable.
\end{frame}

% ============================================================
\begin{frame}{Why interval scaling for \texttt{Tree.scale}?}
The old \code{Tree.scale(s)} multiplied internal heights by \code{s} (affine).
For a heterochronous tree with leaves at \code{h\_c > 0} and parent at \code{h\_N},
the move \textbf{throws} when \code{s < h\_c / h\_N}.

\vspace{1em}

This is not an edge case.
Most BEAST analyses with sampling dates run on heterochronous trees.

\vspace{1em}

\textbf{Interval scaling}: multiply each internal node's margin (height above its taller child) by \code{s}.

\begin{itemize}
\item Tip dates preserved by construction.
\item Move is valid for any positive \code{s}; never throws.
\item Sum of margins is exactly \code{s}-equivariant: matches \code{getScalableValue}.
\item For binary trees, the dof count equals the number of margins,
so operator HRs are unchanged.
\end{itemize}
\end{frame}

% ============================================================
\begin{frame}{Affine semantics retained where needed}
Some callers genuinely want ``set my root to exactly this height'' with affine semantics.

\vspace{1em}

\code{Tree.scaleToRootHeight(double targetHeight)} preserves the old affine behaviour as a separate helper, outside the contract.

\begin{itemize}
\item Multiplies internal heights by \code{targetHeight / oldRoot}.
\item May throw \code{IllegalArgumentException} for invalid configurations.
\item Used by \code{StarBeastStartState} init code.
\item Not part of the Scalable contract.
\end{itemize}

\vspace{1em}

Two methods, two purposes: \code{scale} for moves, \code{scaleToRootHeight} for setup.
\end{frame}

% ============================================================
\begin{frame}{AMVN now reaches its proposed value}
Before: AMVN proposed a target root height $V$, called \code{tree.scale(V/oldRoot)}.

\vspace{0.5em}

For non-ultrametric trees under affine scaling, the move sometimes \textbf{threw}.
For non-ultrametric trees under interval scaling, the root would end up at some height other than $V$.
Either way, AMVN's empirical covariance bookkeeping was wrong.

\vspace{1em}

After: AMVN parameterises trees by sum-of-margins via \code{getScalableValue} / \code{setScalableValue}.
\code{setScalableValue(V)} reduces to \code{scale(V / sum\_margins)}.

\begin{itemize}
\item Sum of margins equals $V$ exactly, by the contract.
\item AMVN's covariance is now updated against the realised state.
\item Side benefit: fixes a latent missing \texttt{return} in spec AMVN's \code{getValue} for trees.
\end{itemize}
\end{frame}

% ============================================================
\begin{frame}{Test plan}
Three layers of testing in \href{https://github.com/CompEvol/beast3/pull/70}{PR \#70}:

\begin{itemize}
\item \textbf{ScalableContractTest harness}: a reusable helper that asserts the three invariants on any Scalable. Per-class tests for \code{RealScalarParam}, \code{RealVectorParam}, \code{Tree} on ultrametric, heterochronous, and leaf-intruding fixtures.

\item \textbf{408 existing beast-base unit tests} all pass.

\item \textbf{Heterochronous MCMC integration}: 10 \code{TipTimeTest} chains under \texttt{-Pslow-tests}.
Real chains on tip-dated data converge to BEAST1 reference values within tolerance.
\textit{Tree height: $15147 \pm 73$, expected $15000 \pm 70$, ESS $\sim$22000.}
\end{itemize}

\vspace{1em}

Statistical-tolerance fix on two slow tests (\code{UpDownOperatorTest}, \code{RealRandomWalkOperatorTest}): tolerances widened to 3 SE using ESS rather than nominal sample count, so they fail less often.
\end{frame}

% ============================================================
\begin{frame}{Downstream impact}
Binding API change.
Any package that implements \code{Scalable} or stores the int return needs to update.

\vspace{1em}

\begin{tabular}{lll}
\toprule
\textbf{Package} & \textbf{Impact} & \textbf{Status} \\
\midrule
Mascot          & none (no \code{Scalable}) & compiles clean \\
beast-classic   & none                       & compiles clean \\
LPhyBeast core  & none                       & compiles clean \\
BEASTLabs       & 2 files                    & migrated, all 83 tests pass \\
\bottomrule
\end{tabular}

\vspace{1em}

BEASTLabs migration in parallel \href{https://github.com/BEAST2-Dev/BEASTLabs/pull/29}{PR \#29}:

\begin{itemize}
\item \code{PrevalenceList.scale} returns log Jacobian; new \code{getScalableValue} returns sum of node times.
\item \code{TreeScaleOperator} consumes log Jacobian directly; HR factored as \code{treeLogJ $-$ 2 log(s)}.
\end{itemize}
\end{frame}

% ============================================================
\begin{frame}{Out of scope (deferred)}
Operator-design questions, separable from the interface contract:

\begin{itemize}
\item Spec \code{UpDownOperator}'s \code{actualScaler = lengthAfter / lengthBefore} logic for tree + up/down combinations.
\item Whether \code{IntervalScaleOperator} should collapse into \code{UpDownOperator} now that \code{Tree.scale} does interval scaling.
\item AMVN's covariance behaviour for non-ultrametric trees beyond the proposal-target fix.
\item Whether tree-aware operators should use the kernel \code{s} or an effective tree-dilation factor for non-tree up/down parameters.
\end{itemize}

\vspace{1em}

The interface contract enables these to be tackled separately on their own merits.
\end{frame}

% ============================================================
\begin{frame}{Timeline}
\begin{itemize}
\item \textbf{2 May}: \href{https://github.com/CompEvol/beast3/pull/70}{PR \#70} and \href{https://github.com/BEAST2-Dev/BEASTLabs/pull/29}{BEASTLabs \#29} ready for review.
\item \textbf{Awaiting}: Remco's review (he owns beast-base).
\item \textbf{Proposed at next stakeholders meeting}: this is the last substantive API change to the beast core before an official release.
Aim for release on the BEAST website mid-June.
\end{itemize}
\end{frame}

% ============================================================
\begin{frame}{Summary}
\begin{block}{The contract}
\code{scale(s)} returns the log Jacobian determinant of a move along a single dilation axis.
\code{getScalableValue} reads position on that axis.
\code{setScalableValue(V) = scale(V / getScalableValue())}.
\end{block}

\vspace{0.5em}

\begin{block}{The Tree side}
\code{Tree.scale} does interval scaling and never throws.
Sum of margins is the dilation summary.
\code{Tree.scaleToRootHeight} preserves affine semantics where they were genuinely wanted.
\end{block}

\vspace{0.5em}

\begin{block}{The result}
Joelle's custom trees can implement the contract directly.
AMVN works correctly on non-ultrametric trees.
408 + 9 + 10 tests pass.
Mid-June beast3 release on track.
\end{block}
\end{frame}

\end{document}
